\subsection{Separately concave residuals and minimum cuts}\label{sec:cuts}

In this subsection $F(x)=\frac12x^{\T}Hx+b^{\T}x+c$ is a rational quadratic,
the retained set $\mathcal K$ is a \emph{core} of $k$ continuous
coordinates, and the recourse set $\mathcal R$ is called the
\emph{residual}. Assume
\begin{equation}\label{eq:cr-class}
\text{(R1)}\ H_{ii}\le0\ (i\in\mathcal R),\qquad
\text{(R2)}\ o_io_jH_{ij}\le0\ (i\ne j,\ i,j\in\mathcal R)\ \text{for some }
o\in\{-1,1\}^{\mathcal R}.
\end{equation}
Nothing is assumed about the entries of $H$ in core rows, and the residual
interaction graph may be dense. A residual coordinate may be continuous or
integer; integer bounds are rounded inward first, and an empty domain is
rejected. Signs $o$ as in (R2) are found, or shown not to exist, by the
following lemma applied to $H_{\mathcal R\mathcal R}$.

\begin{lemma}[Recognizing balance]\label{lem:balance}
For a symmetric $m\times m$ matrix $H$, signs $o\in\{-1,1\}^m$ with
$o_io_jH_{ij}\le0$ for all $i\ne j$ can be found, or shown not to exist, in
time linear in $m$ plus the number of nonzero off-diagonal entries.
\end{lemma}

\begin{proof}
The conditions are $o_i=o_j$ when $H_{ij}<0$ and $o_i=-o_j$ when
$H_{ij}>0$. Breadth-first search fixes one sign in each connected component
of the graph of nonzero off-diagonal entries, propagates the conditions
along a spanning tree and checks every other edge. If an edge violates its
condition, it closes a cycle with an odd number of positive entries, around
which the conditions are inconsistent \cite{Harary1953}; otherwise all
conditions hold.
\end{proof}

\begin{lemma}[Endpoint reduction]\label{lem:cr-endpoint}
Fix $v\in\bar X_{\mathcal K}$ and a box $\prod_{i\in\mathcal R}[\ell_i,u_i]$
whose bounds are integers for integer coordinates. Under (R1), the minimum
of $F(v,\cdot)$ over this continuous box is attained at a point with
$y_i\in\{\ell_i,u_i\}$ for all $i\in\mathcal R$, and it equals the minimum
over the mixed residual domain.
\end{lemma}

\begin{proof}
Start from a minimizer, which exists by compactness. For each residual
coordinate in turn, $F$ as a function of that coordinate alone is
$\frac12H_{ii}t^2$ plus an affine function of $t$, which is concave by
(R1); replace the coordinate by an endpoint with no larger value. The
result is a minimizer whose residual coordinates are endpoints. Its integer
coordinates are integral, so it lies in the mixed domain, which is
contained in the continuous box.
\end{proof}

For $i\in\mathcal R$ put $\alpha_i=\ell_i$, $\delta_i=u_i-\ell_i$ if $o_i=1$,
and $\alpha_i=u_i$, $\delta_i=-(u_i-\ell_i)$ if $o_i=-1$. For a
\emph{label} $z\in\{0,1\}^{\mathcal R}$ let $y(z)$ have coordinates
$y(z)_i=\alpha_i+\delta_iz_i$; as $z$ ranges over $\{0,1\}^{\mathcal R}$,
$y(z)$ ranges over all residual endpoint vectors. In a
network with source $\mathsf s$ and sink $\mathsf t$, $\operatorname{cap}(S)$
denotes the total capacity of the arcs leaving a node set $S$ that contains
$\mathsf s$ but not $\mathsf t$.

\begin{proposition}[Signed cut representation]\label{prop:cr-cut}
(a) Let $\Phi(z)=\Phi_0+\sum_i\mu_iz_i+\sum_{i<j}\omega_{ij}z_iz_j$ on
$\{0,1\}^m$ with $\omega_{ij}\le0$ for all $i<j$, and put
$\omega_{ji}=\omega_{ij}$, $\rho_i=\mu_i+\frac12\sum_{j\ne i}\omega_{ij}$ and
$r_{ij}=-\omega_{ij}/2\ge0$. Build a network on $\{1,\dots,m,\mathsf s,\mathsf t\}$
with arcs $i\to\mathsf t$ of capacity $\max\{\rho_i,0\}$ and
$\mathsf s\to i$ of capacity $\max\{-\rho_i,0\}$ for each $i$, and arcs
$i\to j$ and $j\to i$ of capacity $r_{ij}$ for each pair with
$\omega_{ij}\ne0$. With $S_z=\{\mathsf s\}\cup\{i:z_i=1\}$,
\begin{equation}\label{eq:cr-cutid}
\Phi(z)=\Phi_0+\sum_i\min\{0,\rho_i\}+\operatorname{cap}(S_z)
\qquad(z\in\{0,1\}^m).
\end{equation}
(b) Under~\eqref{eq:cr-class} and for a rational core point $v$, the
function $z\mapsto F(v,y(z))$ on $\{0,1\}^{\mathcal R}$ has the form of (a)
with rational coefficients obtained by exact evaluation of $F$, and
$\omega_{ij}=H_{ij}\delta_i\delta_j\le0$.
\end{proposition}

\begin{proof}
(a) For binary $z$,
$\omega_{ij}z_iz_j=\frac{\omega_{ij}}2(z_i+z_j)-\frac{\omega_{ij}}2\abs{z_i-z_j}$,
so $\Phi(z)=\Phi_0+\sum_i\rho_iz_i+\sum_{i<j}r_{ij}\abs{z_i-z_j}$. In the
cut $S_z$, an arc $i\to\mathsf t$ is cut exactly when $z_i=1$, an arc
$\mathsf s\to i$ exactly when $z_i=0$, and exactly one of the arcs
$i\to j$, $j\to i$ when $z_i\ne z_j$, none otherwise. Hence
$\operatorname{cap}(S_z)=\sum_i[\max\{\rho_i,0\}z_i+\max\{-\rho_i,0\}(1-z_i)]
+\sum_{i<j}r_{ij}\abs{z_i-z_j}$, and
$\rho_iz_i=\max\{\rho_i,0\}z_i+\max\{-\rho_i,0\}(1-z_i)+\min\{0,\rho_i\}$.

(b) Expand $F(v,\alpha+\delta\circ z)$ and use $z_i^2=z_i$. Terms without
residual coordinates are constant. The terms $\frac12H_{ii}y_i^2$, $b_iy_i$
and $H_{ij}v_jy_i$ with $j\in\mathcal K$ are affine in $z_i$, and
$H_{ij}y_iy_j$ with $i\ne j$ in $\mathcal R$ gives
$H_{ij}\delta_i\delta_jz_iz_j$ plus affine terms. By (R2),
$\omega_{ij}=H_{ij}o_io_j(u_i-\ell_i)(u_j-\ell_j)\le0$.
\end{proof}

Every node set that contains $\mathsf s$ but not $\mathsf t$ is $S_z$ for
exactly one $z$, so by~\eqref{eq:cr-cutid} minimizing $\Phi$ is a minimum
$\mathsf s$--$\mathsf t$ cut problem with nonnegative capacities.

\begin{theorem}[Certified cut recourse]\label{thm:cr-oracle}
Assume~\eqref{eq:cr-class}. For every rational core point $v$, one
maximum-flow computation returns a minimizer $y(v)$ of $F(v,\cdot)$ over
$X_{\mathcal R}$, which is a residual endpoint vector, the exact value
$V(v)$, and a certificate: a label $z$ with $y(v)=y(z)$ and a feasible flow
in the network of Proposition~\ref{prop:cr-cut} whose value equals
$\operatorname{cap}(S_z)$. The bit complexity is polynomial in the encoding
length of $F$ and $v$, with an absolute exponent. Given the network, the
certificate is checked with a number of rational operations linear in
$\abs{\mathcal R}$ plus the number of nonzero entries $H_{ij}$ with
$i\ne j$ in $\mathcal R$. The same holds after the residual box is replaced
by a rational subbox, after residual coordinates are fixed, and after
rational linear terms are added.
\end{theorem}

\begin{proof}
By Lemma~\ref{lem:cr-endpoint} the conditional minimum is attained at some
$y(z)$, and by Proposition~\ref{prop:cr-cut} minimizing over $z$ is a
minimum cut problem. Every feasible flow has value at most the capacity of
every cut (weak duality), so a flow of value $\operatorname{cap}(S_z)$
proves that $S_z$ is a minimum cut; checking it needs the capacity bounds,
conservation at each residual node and one comparison. A maximum flow and a
minimum cut of equal value exist and are found, for instance, by the
Edmonds--Karp algorithm \cite{EdmondsKarp1972}, whose number of arithmetic
operations is polynomial in the number of nodes and arcs, here
$\abs{\mathcal R}+2$ and at most $2\abs{\mathcal R}$ plus twice the number
of nonzero residual couplings. After multiplication by a common denominator
of polynomial bit length, the capacities are integers, every augmentation is
by an integer, and every flow value is an integer at most the total
capacity, so all intermediate numbers have polynomial bit length. The
coefficients of Proposition~\ref{prop:cr-cut}(b) are obtained by exact
rational evaluation. The last statement holds because~\eqref{eq:cr-class}
involves only the entries $H_{ij}$ with $i,j\in\mathcal R$, which these
operations do not change; a fixed coordinate has $\delta_i=0$.
\end{proof}

\begin{remark}[Scope and prior work]\label{rem:cr-cut-prior}
Lemmas~\ref{lem:balance} and~\ref{lem:cr-endpoint} and
Proposition~\ref{prop:cr-cut} are classical (Section~\ref{sec:related}).
With an empty core, Theorem~\ref{thm:cr-oracle} is the polynomial
solvability of box QPs that are submodular after sign changes and have
nonpositive diagonal \cite{Padberg1989,BurerNatarajanWillemsen2026v3}. What
we use is the conditional form: the core and the core--residual couplings
are arbitrary, and each conditional value comes with a flow certificate.
Del Pia and Khajavirad also fix the coordinates with nonpositive diagonal at
endpoints; their polynomial classes \cite[Theorem~4]{DelPiaKhajavirad2026}
restrict the structure of these coordinates (logarithmic treewidth and
logarithmic interfaces to the others) and the rank of their coupling to the
others, whereas (R2) restricts only signs and allows a dense residual.
Condition (R1) is what reduces the residual to endpoint vectors. Without it,
and with continuous residual coordinates, the conditional problem is a box
QP that is submodular after sign changes, and the complexity of that class
is open \cite[footnote~2]{BurerNatarajanWillemsen2026v3};
Section~\ref{sec:balanced} grids such coordinates instead
(Corollary~\ref{cor:core}).
\end{remark}

\paragraph{Search over a continuous core.}
Let the core domain be $X_{\mathcal K}=[0,1]^k$ with $k\ge1$, and let the
residual domain $X_{\mathcal R}$ be any compact mixed box. Let $L\ge0$ be an
upper coordinate curvature of $F$ for every core coordinate (for a
quadratic, any $L\ge\max_{i\in\mathcal K}H_{ii}$); by
Lemma~\ref{lem:valuefunction}(a) it is one for $V$. An \emph{exact oracle}
returns, for every dyadic $v\in[0,1]^k$, the value $V(v)$ and a completion
$y(v)\in X_{\mathcal R}$ with $F(v,y(v))=V(v)$, together with a certificate.
Theorem~\ref{thm:cr-oracle} supplies one under~\eqref{eq:cr-class}, and
Remark~\ref{rem:cr-mixed} supplies one when the residual has a convex part.
Level $j$ of the search uses dyadic cells of side $h_j=2^{-j}$, the common
mesh of Section~\ref{sec:growth} with $s=1$, and
\[
e_j=\frac{kLh_j^2}8,\qquad
\beta_V(C)=\min\{V(v):v\text{ a corner of }C\}-e_j .
\]

\begin{algorithm}[CORE: core search]\label{alg:core}
Given an exact oracle, $L$, an accuracy $\varepsilon>0$ and a level limit
$J$, start with $\mathcal Q_0=\{[0,1]^k\}$ and $U=+\infty$. For
$j=0,1,\dots,J$:
\begin{enumerate}[label=(\roman*)]
\item query the oracle at every corner of a cell of $\mathcal Q_j$ that has
not been queried before, and let $U$ be the least value $V(v)$ of all
queried corners, attained by the incumbent $(v_U,y(v_U))$;
\item put $\lambda_j=\min_{C\in\mathcal Q_j}\beta_V(C)$ and retain
$\mathcal Q_j'=\{C\in\mathcal Q_j:\beta_V(C)\le U\}$;
\item if $U-\min\{\lambda_j,U\}\le\varepsilon$, stop; otherwise let
$\mathcal Q_{j+1}$ consist of the $2^k$ dyadic children of each cell of
$\mathcal Q_j'$.
\end{enumerate}
Return the incumbent, $U$ and the lower bound $\min\{\lambda_j,U\}$.
\end{algorithm}

\begin{theorem}[Core search]\label{thm:cr-search}
At every level $j$ of CORE:
\begin{enumerate}[label=(\alph*)]
\item every cell of $\mathcal Q_j$ that contains $x^*_{\mathcal K}$ for a
global minimizer $x^*$ is retained; in particular $\mathcal Q_j'\ne\emptyset$;
\item $\lambda_j\le\OPT\le U\le\lambda_j+e_j$;
\item every retained cell has a corner $v$ with $V(v)\le\OPT+2e_j$;
\item every $x\in X$ satisfies $F(x)\ge\min\{\lambda_j,U\}$, and this can be
verified from the list of cells, the certified queried values and the
retention decisions, without any optimization.
\end{enumerate}
\end{theorem}

\begin{proof}
Let $C$ be a cell of side $h_j$ and $x\in X$ with $x_{\mathcal K}\in C$.
Lemma~\ref{lem:cr-cell} with $B=\mathcal K$ gives
$F(x)\ge\min\{V(v):v\text{ a corner of }C\}-e_{\mathcal K}(C)\ge\beta_V(C)$,
because the core coordinates are continuous and $e_{\mathcal K}(C)\le e_j$.

(a) The unit cube contains $x^*_{\mathcal K}$. If $C\in\mathcal Q_j$ contains
it, then $\beta_V(C)\le F(x^*)=\OPT\le U$, so $C$ is retained and every
child of $C$ that contains $x^*_{\mathcal K}$ belongs to $\mathcal Q_{j+1}$.
(b), (c) These are Theorem~\ref{thm:cr-filter} with $B=\mathcal K$, the
cells $\mathcal Q_j$, $e=e_j$, $\underline V=V$ and $\varepsilon_{\mathrm{or}}=0$; its hypothesis
on an optimizer cell holds by (a), and $\lambda_j\le\OPT$ by the first
paragraph applied to $x^*$.
(d) Children of retained cells cover them, so the cells of $\mathcal Q_j$
together with the cells discarded at levels $i<j$ cover $[0,1]^k$. If
$x_{\mathcal K}$ lies in a cell of $\mathcal Q_j$, then $F(x)\ge\lambda_j$ by
the first paragraph. If it lies in a cell $C$ discarded at level $i$, then
$F(x)\ge\beta_V(C)>U^{(i)}\ge U$, where $U^{(i)}$ is the incumbent value at
level $i$. Each step uses only recorded values.
\end{proof}

By (b), CORE stops at the latest at the first level $j$ with
$e_j\le\varepsilon$, provided that $J$ is at least this level. Its cost is
controlled by the number of near-optimal corners.

\begin{proposition}[Query count under core growth]\label{prop:cr-growth}
Suppose $V(v)-\OPT\ge g\norm{v-v^*}^2$ for all $v\in[0,1]^k$, some $v^*$
and some $g>0$; by Lemma~\ref{lem:valuefunction}(b) this holds with the same
$g$ if $F$ has point growth with constant $g$. Let
$\kappa_{\mathcal K}=\kappa(L,g)=\max\{1,L/g\}$. Then levels $0,\dots,J$ of
CORE make at most
\[
2^k+8^kJ\bigl(1+\sqrt{k\kappa_{\mathcal K}}\bigr)^k
\]
oracle queries. The constant $g$ is not used by the algorithm.
\end{proposition}

\begin{proof}
Let $N_j$ be the number of points $v\in h_j\Z^k\cap[0,1]^k$ with
$V(v)\le\OPT+2e_j$. Level $0$ makes $2^k$ queries. By
Theorem~\ref{thm:cr-search}(c), every retained level-$j$ cell has a corner
counted in $N_j$, and a point is a corner of at most $2^k$ cells, so
$\abs{\mathcal Q_j'}\le2^kN_j$. Each retained cell has $2^k$ children with
$2^k$ corners each, so level $j+1$ makes at most $8^kN_j$ queries. A point
counted in $N_j$ satisfies $g\norm{v-v^*}^2\le2e_j=kLh_j^2/4$, so each of
its coordinates lies within $\frac{h_j}2\sqrt{k\kappa_{\mathcal K}}$ of that
of $v^*$; an interval of length $h_j\sqrt{k\kappa_{\mathcal K}}$ contains at
most $1+\sqrt{k\kappa_{\mathcal K}}$ points of $h_j\Z$. Hence
$N_j\le(1+\sqrt{k\kappa_{\mathcal K}})^k$; sum over $j<J$.
\end{proof}

The bound is linear in the number of levels, and the residual enters only
the cost of each query. Without growth the count can be exponential in the
number of levels.

\begin{example}[Flat directions]\label{ex:cr-flat}
Take $k$ even, no residual, and
$V(v)=\sum_{t=1}^{k/2}(v_{2t-1}-v_{2t})^2$, so $L=2$ and $\OPT=0$. The corner $v=0$ gives $U=0$ at level $0$. Every product
of diagonal cells
$\prod_t\bigl([a_th_j,(a_t+1)h_j]\times[a_th_j,(a_t+1)h_j]\bigr)$ with
integers $0\le a_t<2^j$ has a
corner with $V=0$, hence is retained, and its parent is again such a
product. So level $j$ retains at least $2^{jk/2}$ cells, while
$U-\lambda_j=e_j$; stopping at accuracy $\varepsilon$ therefore takes
$\Omega(\varepsilon^{-k/4})$ queries for fixed $k$.
\end{example}

\paragraph{Exact output.}
Under~\eqref{eq:cr-class} and for rational data, $\OPT$ and every value $\min_{v\in[0,1]^k}F(v,y)$ at a
residual endpoint vector $y$ have denominators at most a constant
$\Omega_{\mathcal K}$ of bit length $O(I^2)$ (Lemma~\ref{lem:cr-height}).
Run CORE with the oracle of Theorem~\ref{thm:cr-oracle} and accuracy
$\varepsilon=\Omega_{\mathcal K}^{-2}/2$, and minimize $F(\cdot,y)$
exactly over the faces of the core cube, where $y$ is the residual part of
the incumbent. The result is a global minimizer that passes the acceptance
test of Proposition~\ref{prop:accept} with the denominator bound
$\Omega_{\mathcal K}$. Under core growth the total cost is
$8^k(1+\sqrt{k\kappa_{\mathcal K}})^k\poly(I)$ bit operations with an
absolute exponent, and $g$ is not an input (Theorem~\ref{thm:cr-exact} in
Appendix~\ref{app:recourse-cuts}).
